\documentclass{report}

\def\filename{Tut1Q3}
\def\course{UECM3393}
\def\chapsec{Tutorial 1}
\def\date{\today} % to be changed

\def\headL{\textbf{\course :} \chapsec}
\def\headC{}
\def\headR{\date}

\def\footL{}
\def\footC{\thepage}
\def\footR{Carlos Wong 2303326}

\usepackage{amsmath, amssymb, amsfonts}

\usepackage[a4paper, hmargin=0.85in, vmargin=0.9in, includeheadfoot]{geometry}

\usepackage{fancyhdr}
\setlength{\headheight}{14pt}
\pagestyle{fancy}
\fancyhf{}

\lhead{{\large \headL}}
\chead{{\large \headC}}
\rhead{{\large \headR}}

\lfoot{\footL}
\cfoot{\footC}
\rfoot{\footR}

\usepackage{etoolbox}



\begin{document} 
\paragraph{Q3. $\\$}

(a) Show that if any 14 integers are selected from the set $S = \{1,2,3, \dots , 25\}$, there are at least two whose sum is 26.

\medskip \noindent \textit{Solution.} 

We see that $1+25=26, 2+24=26, \dots , 12+14=26$.

That is, there are 12 distinct pairs of integers $(a, b)$ that can be chosen from the set $S$ s.t. $a+b=26$.

Note that if 13 is selected, it has no corresponding pair such that they sum to 26.

From here, we can construct 13 pigeonholes from the 12 distinct pairs of integers and the one lone integer that has no corresponding pair.

By the Pigeonhole Principle, putting 14 pigeons into 13 pigeonholes guarantees that one pigeonhole has at least two pigeons.

Thus, selecting 14 integers from $S$ guarantees one pair of integers that sum to 26 being chosen. \hfill $\square$

\bigskip

(b) Write a statement that generalizes the results of part (a).

\medskip \noindent \textit{Solution.}

Selecting $\lfloor \frac{n+1}{2}\rfloor + 1$ integers from $S = \{1,2, \dots , n\}$ guarrantees that there are at least two integers selected that sum to $n+1$. \hfill $\square$

\end{document}
